\scanpage{73}
% archive sheet number: 72
% margin note written sideways
\margin{Petits fourbis \unsure{chomologiques|cohomologiques}}
% L1-9
\uline{Proposition} Soit $k$ un corps commutatif, $\unsure{L/k}$ une extension galoisienne de $k$ de groupe de Galois $\Gamma$, $\underline A$ une algèbre avec 1 de dimension finie sur $k$, $\underline A^L$ l'algèbre sur $L$ obtenue par extension des scalaires à $L$, $A^L$ est considéré comme muni du groupe d'automorphismes d'anneau $\Gamma$. Soit $G=\underline A^{L*}$ le groupe des éléments inversibles de $\underline A^L$, $G$ \deleted{\unsure{$\underline A^L$|est}} un $\Gamma$-groupe.

$H^1(\Gamma,G)$ est réduit à l'élément neutre, ainsi que $H^1(\Gamma,G^\psi)$ si $\psi$ est un hom croisé de $\Gamma$ dans $\deleted{\operatorname{Int}G}\operatorname{Aut}\underline A$.

% L10
\uline{Corollaire} $H^1(\Gamma,\operatorname{Gl}(n,L))$ est réduit à l'élément neutre.

% L11-12
\uline{Démonstration} $\Gamma$ opère dans $G^\psi=G$

$\sigma(\gamma)g=\psi(\gamma).g^\gamma$. Soit $\unsure{\varphi|\phi}\in Z^1(\Gamma,G^\psi)$, \deleted{\unsure{posons}} faisons
% L13
\[
\deleted{\unsure{a}}\quad x=\sum_{\gamma\in\Gamma}\bigl(\deleted{\unsure{\sigma(\gamma).A}}\bigr)\varphi(\gamma)\sigma(\gamma).A\qquad(A\in\underline A^L)
\]
% L14, cancelled line
\deleted{les \unsure{$\varphi(\gamma)$}\ \unsure{les|et les}\ $\sigma(\gamma)$ étant \unsure{en fait}\ \illegible{1}}
\deleted{\unsure{$\exists$}}\deleted{\illegible{1}}
% L15-18
\[
\begin{aligned}
\sigma(\gamma)x&=\sum_{\gamma'\in\Gamma}\bigl(\deleted{\unsure{\sigma(\gamma\gamma')A}}\bigr)\sigma(\gamma)\varphi(\gamma')\sigma(\gamma\gamma').A\\
&\deleted{=\sum_{\gamma'\in\Gamma}(\sigma(\gamma')A)\sigma(\gamma)\varphi(\gamma^{-1}\gamma')}\\
&=\sum_{\gamma'\in\Gamma}\sigma(\gamma)\varphi(\gamma^{-1}\gamma')\sigma(\gamma')A\\
&=\sigma(\unsure{\gamma})\varphi(\gamma^{-1})\sum_{\gamma'\in\Gamma}\varphi(\gamma')\sigma(\gamma')A=\sigma(\gamma)\varphi(\gamma^{-1})x
\end{aligned}
\]
% L19
\[
x=\varphi(\gamma^{-1})\sigma(\gamma)^{-1}x\qquad x=\varphi(\gamma)\sigma(\gamma)x
\]
% L20-21
$\deleted{\unsure{\varphi(\gamma)=}}$ D'où, \uline{si $x$ inversible}, la formule
\[
\varphi(\gamma)=x\bigl(\sigma(\gamma)x\bigr)^{-1}
\]
% L22-23
Il reste à prouver donc que l'on peut trouver $A$ de telle façon que $\sum\varphi(\gamma)\sigma(\gamma).A$ soit
